\documentclass[10pt,a4paper]{amsart}
\usepackage[margin=19mm]{geometry}
\usepackage{amsmath,amssymb,mathtools}
\usepackage[scr=rsfs,cal=euler]{mathalfa}
\usepackage{xfrac}
\usepackage[unicode,hidelinks,bookmarks=false]{hyperref}
\newcommand{\ie}{i.e.\ }
\newcommand{\cf}{cf.\ }
\newcommand{\resp}{respectively\ }
\newcommand{\Subsection}{\subsection}
\providecommand{\nobreakdash}{\nobreak\hbox{-}}
\setlength{\emergencystretch}{2em}
\multlinegap0pt
\usepackage{bm}
\usepackage{bbm}
\usepackage{dsfont}
\usepackage{xspace}
\usepackage{cancel}
\usepackage{mathtools}
\usepackage{amsthm}
\makeatletter
\@ifundefined{theorem}{\newtheorem{theorem}{Theorem}}{}
\makeatother
\title[Analytic Dilogarithm Identities]{Analytic Dilogarithm Identities}
\author{Cetin Hakimoglu-Brown}
\date{Source: arXiv:2506.10206v2 (CC BY 4.0)}
\begin{document}
\maketitle

\noindent\emph{Source and licence.} Analytic Dilogarithm Identities. Version arXiv:2506.10206v2 (CC BY 4.0), distributed under the Creative Commons Attribution 4.0 International licence, as recorded in the arXiv licence metadata for this exact version and shown on the abstract page https://arxiv.org/abs/2506.10206v2. Full rights notice: licenses/arxiv-2506.10206-CC-BY-4.0.txt.\par\smallskip

\noindent\textit{Editorial scope.} Counted unit: the Theorem whose source label is \texttt{theoremSunCampbell} in the article (source0.tex lines 362--367, source label \texttt{theoremSunCampbell}), delivered together with the article's own complete original proof of it (source0.tex lines 369--408, one \texttt{proof} environment of 40 lines). No mathematics is written, repaired or re-derived: statement and proof are the article's own text line for line. Required context retained uncounted: maintheorem (lines 262--316); mainABCD (lines 297--299); MequalsAB (lines 301--305); displays3 (lines 307--312). Every definition, display and label of those lines is retained unchanged. Omitted from this package and not redistributed: 1--261, 300--300, 306--306, 313--361, 368--368, 409--2050. Nothing inside a retained excerpt is omitted or altered: every retained body is delivered exactly as the article's lines are written, so no omission receipt of any kind accompanies this package. Citations: the retained excerpts carry no citation command at all, so no bibliography entry of the article is transcribed and no reference is redistributed. Reference adaptation: none was needed and none was made; every reference inside the delivered unit resolves inside it, so no unresolved reference or citation is produced. Printed slips kept literal: no printed word, symbol, hypothesis or number of the counted statement or of its proof was altered. Layout and class: the article's own class is \texttt{article} with packages including amssymb, amsmath, amsthm, fullpage, xcolor; the wrapper is the lane's pinned \texttt{amsart} editorial wrapper at 10pt on A4, which restates the theorem-like environments the retained text needs and adds the article's own macro definitions that the retained text needs (none), with extra wrapper packages (bm, bbm, dsfont, xspace, cancel, mathtools). The delivered numbering is the unit's own. Assets: no retained span contains a figure environment, a graphics-inclusion command, a PDF-page inclusion, a table or a TikZ picture; no image file is copied into this package, the package declares no asset dependency and no image file is redistributed with it. Licence: version arXiv:2506.10206v2 of this article is distributed under the Creative Commons Attribution 4.0 International licence, as recorded in the arXiv licence metadata for this exact version and shown on the abstract page https://arxiv.org/abs/2506.10206v2. A full rights notice is delivered at licenses/arxiv-2506.10206-CC-BY-4.0.txt. Characters: every retained excerpt is pure ASCII, so the delivered engine, XeLaTeX with its default Latin Modern text font, reports no missing character.
\begin{theorem}\label{maintheorem}
 Let 
\begin{align*}
 A & := \operatorname{Li}_2\left(\frac{\sqrt{4 - 3u^2} + u^2 - 2}{u (u + 1)}\right) 
 - \operatorname{Li}_2\left(\frac{\sqrt{4 - 3u^2} + u^2 - 2}{(u - 1) u}\right), \\ 
 B & := 
 \frac{1}{2} \ln^2\left(\frac{-i \sqrt{3 u^2-4}+u+2}{4} \right) - 
 \frac{1}{2} \ln ^2\left(\frac{- i \sqrt{3 u^2 - 4} - u^2 +2}{u 
 (u+1)}\right) + \\ 
 & \quad \frac{1}{2} \ln ^2\left(\frac{i \sqrt{3 u^2 - 4} + u + 2}{4} \right)-\frac{\pi ^2}{6}, \\ 
 J & := \frac{1}{4} \ln^2 \left( \frac{-4}{-i \sqrt{3 u^2 - 4} + u + 2} \right) + 
 \frac{1}{4} \ln^2 \left( \frac{-4}{i \sqrt{3 u^2 - 4} + u + 2} \right) - \\ 
 &\quad \frac{1}{4} \ln^2 \left( \frac{-4}{-i \sqrt{3 u^2 - 4} - u + 2} \right) - \frac{1}{4} 
 \ln^2 \left( \frac{-4}{i \sqrt{3 u^2 - 4} - u + 2} \right), \\ 
 C & := - \frac{1}{2} \ln \left( \frac{ - i \sqrt{3 u^2 - 4} + u + 
 2}{4} \right) \ln \left( \frac{ i \sqrt{3 u^2 - 4} - u + 2}{4} \right) - \\ 
 & \quad \frac{1}{2} \ln \left( \frac{i \sqrt{3 u^2 - 4} + u + 2}{4} \right) 
 \ln \left( \frac{- i \sqrt{3 u^2 - 4} - u + 2}{4} \right) + \frac{\pi^2}{6}, \\ 
 H & := 
 \frac{1}{4} \ln^2\left(\frac{2}{u-1}\right)-\frac{1}{4} \ln ^2\left(-\frac{2}{u+1}\right)+\frac{1}{2} \ln \left(\frac{-u + 1}{2} \right) 
 \ln \left(\frac{u+1}{2}\right) - \\ 
 & \quad \frac{\pi i}{2} \ln \left(\frac{u + 
 1}{u-1}\right)-\ln (2) \ln \left(\frac{u+1}{u-1}\right)-\frac{\pi^2}{12}, \\ 
 D & := \operatorname{Li}_2 \left( \frac{1 - u}{2} \right) 
 + \ln(2) \ln \left( \frac{u - 1}{u + 1} \right) + \left( 2 \ln(2) + \frac{\pi i}{2} \right) \ln \left( \frac{u + 
 1}{u - 1} \right), \\ 
 K & := \frac{1}{2} \operatorname{Li}_2 \left( \frac{2}{- \sqrt{4 - 3 u^2} + u} \right) - 
 \frac{1}{2} \operatorname{Li}_2 \left( \frac{2}{\sqrt{4 - 3 u^2} - u} \right) \\ 
 &\quad - \frac{1}{2} \operatorname{Li}_2 \left( \frac{-2}{\sqrt{4 - 3 u^2} + u} \right) + 
 \frac{1}{2} \operatorname{Li}_2 \left( \frac{2}{\sqrt{4 - 3 u^2} + u} \right) + 
 \frac{1}{2} \operatorname{Li}_2 \left( \frac{-1}{u} \right) - 
 \frac{1}{2} \operatorname{Li}_2 \left( \frac{1}{u} \right). 
\end{align*} 
 Then, for all $u \in \mathbb{C} \setminus \mathbb{R}$ such that $\Re u \geq 0$ and such that the power series 
 in \eqref{displays3} converges, we have that 
\begin{equation}\label{mainABCD}
 A + B + J + C + H + D = s_3 = K, 
\end{equation}
 and that 
\begin{equation}\label{MequalsAB}
 M = A + B = 
 - \operatorname{Li}_{2}\left( \frac{\sqrt{4 - 3 u^2} - u + 2}{4} \right) 
 - \operatorname{Li}_{2}\left( \frac{-\sqrt{4 - 3 u^2} - u + 2}{4} \right), 
\end{equation}
 writing 
\begin{equation}\label{displays3}
 s_3 = \frac{2}{3u(u^2-1)} \sum_{n=0}^{\infty} \left( \frac{1}{u(u^2-1)} \right)^{2n} \left( \frac{4}{27} \right)^{n} 
 \left[ \begin{matrix} \frac{1}{2}, \frac{1}{2}, 1 
 \vspace{1mm} \\ 
 \frac{5}{6}, \frac{7}{6}, \frac{3}{2} \end{matrix} \right]_{n}. 
\end{equation}
 Moreover, for $u \in \mathbb{R}_{\geq 0}$ such that \eqref{displays3} converges, by letting $v \in \mathbb{C} \setminus \mathbb{R}$ 
 approach $u$ and taking the corresponding limits on all sides of the equalities among \eqref{mainABCD}, these limits are equal, 
 and similarly for \eqref{MequalsAB}. 
\end{theorem}

\begin{theorem}\label{theoremSunCampbell}
 The closed-form evaluation
 $$ 8 \operatorname{Li}_2\left(2-\sqrt{3}\right)-\operatorname{Li}_2\left(4 \sqrt{3} - 7 \right) 
 = \frac{5 \pi ^2}{12}+\ln \left(2-\sqrt{3}\right) \ln \left(2+\sqrt{3}\right) $$
 holds true.
\end{theorem}

\begin{proof}
 For the $u = \frac{1}{\sqrt{3}}$ case Theorem \ref{maintheorem}, 
 the left-hand side of \eqref{mainABCD} 
 may be expressed in terms of 
 $ -2 \operatorname{Li}_2 \left( 7 - 4 \sqrt{3} \right) + 
 4 \operatorname{Li}_2 \left( \sqrt{3} - 2 \right) $ and a closed-form expression, 
 and this follows from the known relation 
\begin{equation*}
-2 \operatorname{Li}_2 \left( \frac{3 - \sqrt{3}}{6} \right) = 2 \operatorname{Li}_2 \left( \sqrt{3} - 2 \right) + 
 \ln^{2} \left( \frac{3 + \sqrt{3}}{6} \right). 
\end{equation*}
 The right-hand side of \eqref{mainABCD}, for the same input value 
 $u = \frac{1}{\sqrt{3}}$ may be rewritten as 
 $$ K = \operatorname{Li}_2\left(-\sqrt{3}\right)-\frac{1}{2}
 \operatorname{Li}_2\left(-\frac{\sqrt{3}}{2}\right) + 
 \frac{\operatorname{Li}_2\left(\frac{\sqrt{3}}{2}\right)}{2}-\operatorname{Li}_2\left(\sqrt{3}\right). $$
 Using the well known $4$- term functional equation 
\begin{equation}\label{wellknown4term}
 \operatorname{Li}_2(z) - \operatorname{Li}_2(-z) + 
 \operatorname{Li}_2\left(\frac{1 - z}{1 + z}\right) - 
 \operatorname{Li}_2\left(\frac{-(1 - z)}{1 + z}\right) = \frac{\pi^2}{4} + \ln(z) \ln\left(\frac{1 + z}{1 - z}\right), 
\end{equation}
 we obtain, by setting $z = \frac{\sqrt{3}}{2}$ and $z=\sqrt{3}$, that the right-hand side of \eqref{mainABCD} may be rewritten as 
\begin{equation*}
 \operatorname{Li}_2\left(7 - 4 \sqrt{3}\right) - \operatorname{Li}_2\left(4 \sqrt{3} - 7 \right) - 
 2 \, \operatorname{Li}_2\left(\sqrt{3} - 2\right) + 2 \, \operatorname{Li}_2\left(2 - \sqrt{3}\right). 
\end{equation*}
 Equating our equivalent expressions for the right-hand and left-hand sides of \eqref{mainABCD}, we obtain 
\begin{equation*}
 3 \, \operatorname{Li}_2\left(7 - 4 \sqrt{3}\right) - 
 \operatorname{Li}_2\left(4 \sqrt{3} - 7\right) - 
 6 \, \operatorname{Li}_2\left(\sqrt{3} - 2\right) + 2 \, \operatorname{Li}_2\left(2 - \sqrt{3}\right)
\end{equation*}
 and 
\begin{equation*}
 3 \, \operatorname{Li}_2\left(7 - 4 \sqrt{3}\right) = 6 \, 
 \operatorname{Li}_2\left(2 - \sqrt{3}\right) + 6 \, \operatorname{Li}_2\left(\sqrt{3} - 2\right), 
\end{equation*}
 and this gives us an equivalent version of the desired result. 
\end{proof}

\end{document}
